% ==================================================================
% id:       seq-012
% topic:    סדרות
% subtopic: סדרות מונוטוניות וסדרות נסיגה
% solution: full
% ==================================================================
\textbf{יהיו
$a_1\ge b_1$
מספרים חיוביים.
נגדיר ברקורסיה שתי סדרות באופן הבא:}
\[ .a_{n+1} = \frac{a_n+b_n}{2},~ b_{n+1} = \sqrt{a_nb_n}\]
\begin{enumerate}
\item
הוכיחו שלכל
$n \in\nn$
מתקיים
$b_{n} \le b_{n+1} \le a_{n+1} \le a_{n}$.
\hint{השתמשו באינדוקציה ובאי-שוויון הממוצעים $\sqrt{xy}\le\frac{x+y}{2}$ לכל $x,y>0$.}
\begin{solbox}
ניעזר באינדוקציה ובאי-שיוויון הממוצעים שאומר כי לכל שני מספרים חיוביים
$x,y\in\rr$
מתקיים
$$.\sqrt{xy}\le \frac{x+y}{2}$$
\begin{itemize}
     \item
עבור
$n=1$
מראים כי אכן
$$b_1\le b_2\le a_2\le a_1$$
מה שנובע מההנחה כי
$b_1\le a_1$
ומאי-שיוויון הממוצעים (בדקו).
\item
נניח כי
$b_n\le b_{n+1}\le a_{n+1}\le a_n$
ומכך נסיק, שוב בעזרת אי שיוויון הממוצעים ובעזרת הנחת האינדוקציה, כי
$$b_{n+1}\le b_{n+2}\le a_{n+2}\le a_{n+1}$$
בדקו זאת!
 \end{itemize}
\end{solbox}
\item
הראו כי הסדרות
$(a_n)$
ו-
$(b_n)$
מתכנסות.
\begin{solbox}
    מהסעיף הראשון נובע כי שתי הסדרות הן מונוטוניות וחסומות, לכן מתכנסות.

\end{solbox}
\item
הראו שמתקיים
$.\lim_{n \to \infty}a_n
= \lim_{n \to \infty}b_n$

הגבול המשותף נקרא
הממוצע הגאומטרי-אריתמטי
של
$(a_n)$
ו-
$(b_n)$.
\hint{השאיפו את $n$ לאינסוף בשני האגפים של נוסחת הנסיגה של $a_{n+1}$.}
\begin{solbox}
    נסמן
$a_n\goes L_1$
ו- $b_n\goes L_2$.
משיקולי אריתמטיקה נקבל
$$\frac{a_n+b_n}{2}\rightarrow\frac{L_1+L_2}{2}$$
כאשר מצד שני
$a_{n+1}\goes L_1$,
לכן מיחידות הגבול (של סדרה מתכנסת) נקבל כי
$$\frac{L_1+L_2}{2}=L_1$$
ואז
$L_1=L_2$.
\end{solbox}
\end{enumerate}
